\section{Salinité des sols et stratégies de correction}
\subsection{Problématique de la salinité des sols dans la zone d'étude}
La salinité des sols est l’un des principaux facteurs de dégradation des terres dans le Sine-Saloum, particulièrement dans la région de Fatick. Selon Amar \cite{noauthor_salinisation_nodate}, plus de 1,7 million d'hectares de terres arables au Sénégal sont affectées par le sel, conséquence de l’intrusion marine, de la remontée capillaire des sels et de  la variabilité climatique. Dans les zones côtières, 
la salinité transforme progressivement les rizières et les champs maraîchers en tannes improductives, réduisant ainsi la disponibilité des terres cultivables. Cette situation compromet la sécurité alimentaire et pousse de nombreux producteurs à abandonner leurs terres.


\subsection{Effet du stress salin sur les cultures maraîchères}
Le stress salin agit à plusieurs niveaux sur les cultures maraîchères. D’une part, il réduit la germination des semences et ralentit la croissance des jeunes plants. D’autre part, il provoque une toxicité ionique due à l'accumulation de sodium ($\mathrm{Na^{+}}$) et de chlorure ($\mathrm{Cl^{-}}$), qui perturbe l’absorption des nutriments essentiels comme le potassium et le calcium. Dans les Niayes et le bassin arachidier, les cultures de tomates, oignons et chou présentent des baisses de rendement pouvant dépasser 50\% \cite{fall_salt_nodate}. Les observations de Fall Ba et al. \cite{fall_salt_nodate} montrent que la salinité  peut entraîner une réduction significative de la photosynthèse et est associé à une augmentation du flétrissement prématuré, ce qui fragilise la rentabilité des exploitations maraîchères.

\subsection{Cultures étudiées et sensibilité à la salinité}
Trois cultures maraîchères ont été retenues : la laitue (Lactuca sativa), le navet (Brassica rapa) et le gombo (Abelmoschus esculentus). Le tableau ci-dessous présente les caractéristiques utiles à cette étude.

{\small
\begin{longtable}{|p{2cm}|p{3.2cm}|p{4.4cm}|p{4.4cm}|}
    \caption{Cycle, sensibilité au sel et besoins en eau des trois cultures}
    \label{tab:cultures_sel_eau} \\
    \hline
    \textbf{Culture} & \textbf{Cycle} & \textbf{Sensibilité au sel} & \textbf{Besoins en eau} \\
    \hline
    \endfirsthead
    \hline
    \textbf{Culture} & \textbf{Cycle} & \textbf{Sensibilité au sel} & \textbf{Besoins en eau} \\
    \hline
    \endhead
    Laitue & 40 à 50 jours du semis à la récolte & Sensibilité modérée, seuil de 1,3 à 1,75~dS/m ; effet principal : osmotique & Arrosage régulier et contrôlé (racines très limitées) ; manque : arrêt de croissance ; excès : asphyxie des racines \\
    \hline
    Navet & 45 à 48 jours du semis à la récolte & Moins sensible que la laitue mais moins résistant que le gombo & Irrigation régulière, surtout à l'expansion de la racine ; supporte mal l'alternance sécheresse et saturation \\
    \hline
    Gombo & 2 à 3 mois & Plus résistant que la laitue et le navet & Besoins constants, surtout à la floraison et à la fructification ; un déficit provoque la chute des fleurs et des fruits \\
    \hline
\end{longtable}
}

Ces exigences, en particulier la sensibilité au sel et la régularité des apports en eau, justifient le pilotage de l’irrigation étudié dans ce mémoire.

\subsection{Amendements organiques utilisés pour lutter contre la salinité (biochar, compost, coque d’arachide et biofertilisant)}
Face à ce problème, les amendements organiques apparaissent comme une solution durable et adaptée au contexte sénégalais. Le biochar, obtenu par  pyrolyse de biomasse, améliore la capacité de rétention d'eau et favorise la fixation des sels, réduisant ainsi leur disponibilité pour les plantes. Les coques d’arachides, abondantes dans le bassin arachidier, sont utilisées comme paillage ou compost. Elles contribuent à diminuer la conductivité électrique des sols et à améliorer leur structure. Le compost, riche en MO, stimule l'activité microbienne et augmente la tolérance des cultures au stress salin. Enfin, les fertilisants apportent des  micro-organismes bénéfiques qui favorisent la croissance des plantes en conditions difficiles. Les travaux de Diédhiou et al. \cite{diedhiou_investigation_2023} en Casamance ont montré que l’application combinée de biochar et de compost améliore la survie et le rendement du riz en condition saline, avec des taux de survie supérieurs à 90\%.

\subsection{Travaux antérieurs sur l'efficacité de ces amendements}
Les recherches menées par Niang et al. \cite{niang_degradation_2017} dans le Haut Saloum et le delta du Sénégal ont démontré que l’apport de compost et l’utilisation de CA permettent une réhabilitation progressive des sols salés. Ces pratiques agroécologiques favorisent la restauration de la fertilité et la reprise des cultures maraîchères. De plus, les initiatives locales de l’ISRA et du CIRAD dans le département de Fatick ont confirmé que la combinaison d’amendements organiques et de reboisement avec des espèces halophytes (Casuarina, Melaleuca) constitue une approche durable pour restaurer les sols et améliorer la résilience des systèmes agricoles. \textbf{Ces résultats montrent que les solutions endogènes, adaptées aux réalités locales, sont plus efficaces que les interventions purement chimiques}.


